\subsection{OPERATIONS ON INTEGERS AND FINITE SETS}

PROPOSITION 1. Let \((a_i)_{i\in I}\) be a finite family of integers. Then the cardinals \(\sum_{i\in I}a_i\) and \(\prod_{i\in I}a_i\) are integers.

Let us begin by showing that if \(a\) and \(b\) are integers, then so is \(a + b\) . We proceed by induction on \(b\) . The assertion is true for \(b = 0\) since \(a + 0 = a\) . If \(a + b\) is an integer, then so is \((a + b) + 1\) ( \(\S 4\) , no. 1, Proposition 1). But \((a + b) + 1 = a + (b + 1)\) ( \(\S 3\) , no. 3, Corollary to Proposition 5); hence \(a + (b + 1)\) is an integer, and consequently \(a + b\) is an integer for all integers \(b\) .

Next we show by induction on \(n = \operatorname{Card}(\mathbf{I})\) that \(\sum_{i \in \mathbf{I}} a_i\) is an integer. This is clear if \(n = 0\) , for then \(\mathbf{I} = \emptyset\) and \(\sum_{i \in \mathbf{I}} a_i = 0\) . If \(\operatorname{Card}(\mathbf{I}) = n + 1\) , we may write \(\mathbf{I} = \mathbf{J} \cup \{k\}\) , where \(\operatorname{Card}(\mathbf{J}) = n\) and \(k \notin \mathbf{J}\) . Then

\[
\sum_ {i \in \mathbf {I}} a _ {i} = a _ {k} + \sum_ {i \in \mathbf {J}} a _ {i}
\]

(§ 3, no. 3, Proposition 5). The inductive hypothesis is that \(\sum_{i\in J}a_i\) is an integer; hence so is \(a_{k} + \sum_{i\in J}a_{i}\) , from the first paragraph of the proof. This shows that \(\sum_{i\in I}a_{i}\) is an integer for all \(n\) .

Since the product ab of two integers a and b is the sum of a finite family of integers all equal to a (§3, no. 4, Proposition 6, Corollary 2), ab is an integer. We shall prove by induction on n = Card (I) that \(\prod_{i\in I}a_{i}\) is an integer. This is true for n = 0, because then

\[
\prod_ {i \in \mathbf {I}} a _ {i} = 1.
\]

If Card (I) = n + 1, we have (with the same notation as above)

\[
\prod_ {i \in \mathbf {I}} a _ {i} = a _ {k}. \prod_ {i \in \mathbf {J}} a _ {i}
\]

(§ 3, no. 3, Proposition 5), and the inductive hypothesis therefore implies that \(\prod_{i\in I}a_i\) is an integer. Consequently \(\prod_{i\in I}a_i\) is an integer for all \(n\) .

COROLLARY 1. The union E of a finite family \((\mathbf{X}_i)_{i\in \mathbf{I}}\) of finite sets is a finite set.

For the sum S of the family \((\mathbf{X}_{i})\) is finite; and since there exists a mapping of S onto E (Chapter II, §4, no. 8), the set E is finite (§4, no. 2, Proposition 2, Corollary 3).

COROLLARY 2. The product of a finite family of finite sets is a finite set.

COROLLARY 3. If \(a\) and \(b\) are integers, \(a^b\) is an integer.

For \(a^b\) is the product of a finite family of integers all equal to \(a\) (§ 3, no. 5, Proposition 10).

COROLLARY 4. The set of subsets of a finite set E is finite.

For its cardinal is \(2^{\text{Card (E)}}\) ( \(\S 3\) , no. 5, Proposition 12).

